\chapter{Objectives and Methodology}
\label{chapter:objectives-and-methodology}

Chapters~\ref{chapter:sota} and~\ref{chapter:related} identified two gaps, namely that query-based
\ac{pnp} models as ViP3D \cite{Gu2023ViP3D} create their agent queries from a learned embedding that does not depend on the
scene, and that the consequences of changing this have not been measured in tracking and prediction.
This chapter states the objectives that follow and the methodology by which they are pursued.

\section{Objectives}
\label{sec:objectives}

The goal of this thesis is to develop and evaluate a deep learning architecture for \ac{e2e} perception and
prediction in which the agent queries originate in the point cloud and not in a learned embedding. The
goal is pursued through three objectives.

\begin{itemize}
    \item \textbf{Objective 1.} Design and implement a query initialisation mechanism that derives the 
    agent queries from the \ac{lidar} \ac{bev} feature map, and
    integrate it with the tracker inside ViP3D \cite{Gu2023ViP3D} codebase so that queries carried over from the previous frame keep their slots and
    the initialisation concerns newly appearing objects alone. The resulting model needs to evaluated next to other 
    representants of \ac{e2e} trajectory prediction.

    \item \textbf{Objective 2.} Determine what the camera modality still contributes once the queries already carry
    geometric evidence by implementing multilayer architecture for modality fusion. 

    \item \textbf{Objective 3.} Determine what a query that already begins near its object buys in decoder
    depth and in inference latency emphasizing the trade-off between prediction accuracy and computational cost. 
\end{itemize}

Objective 1 addresses the first gap. Objectives 2 and 3 address the second gap mentioned in \ref{sec:summary}.

\section{Methodology}
\label{sec:methodology}

The work follows the research cycle of analysis, design, implementation, experimental evaluation and
interpretation. Each step is carried by one chapter of this thesis. A brief description of the steps 
is provided in this section.

\begin{itemize}
    \item \textbf{Step 1, analysis.} Survey trajectory prediction and 3D object detection literature, with attention
    to how query-based models initialize their agent queries. The
    outcome is the two gaps and the starting points stated in Section~\ref{sec:summary}. Carried by
    Chapters~\ref{chapter:sota} and~\ref{chapter:related}.

    \item \textbf{Step 2, design.} Propose an architecture that closes those gaps, deciding how a query is
    created from the point cloud, at which point the camera modality re-enters and how the training is
    staged. Designing it mainly concerns designing the initialization mechanisms, design of the fusion decoder layers,
    design of integration with the ViP3D codebase and finally scheduling the training according to the 
    identified literature. Carried by Chapter~\ref{chapter:design}.

    \item \textbf{Step 3, implementation.} Realise the design by extending an existing open implementation
    of ViP3D \cite{Gu2023ViP3D}. Keeping the surrounding code
    fixed is a methodological choice, since it makes the proposed change the only difference between the
    baseline and the proposed model. This way direct comparison can be achieved.
    Carried by Chapter~\ref{chapter:implementation}.

    \item \textbf{Step 4, experimental evaluation.} Measure the implemented model against selected
    \ac{e2e} \ac{pnp} models from the literature, and isolate the contribution of each added component by
    ablating it, one ablation per objective. Detection, tracking and prediction metrics are recorded for
    every configuration and reported per class, together with inference latency and parameter count.
    Chapter~\ref{chapter:experiments-and-evaluation}.

    \item \textbf{Step 5, interpretation.} Given large number of experiments needed to conduct,
    proper interpretation of the results is crucial. The results are analysed using metrics provided in
    Chapter~\ref{sec:sota-metrics} and interpreted in the context of 
    the objectives and the research questions for clear conclusions and recommendation for future work.
    Carried by both Chapter~\ref{chapter:experiments-and-evaluation} and
    Chapter~\ref{chapter:conclusion}.
\end{itemize}

This chapter derived three objectives from the gaps identified in Chapter~\ref{chapter:related} and set
out the research cycle through which they are pursued, from the analysis that produced the gaps to the
interpretation of the results. In the following chapter, the design of the proposed model is described in detail.
